% !TEX root = ../tfg.tex

% IDEAS (borrar antes de pasar la memoria al tutor)
% - Para quién: clubes que quieran este análisis y tengan, o pongan, una cámara
%   fija al fondo de la pista, en alto.
% - Entra: el vídeo crudo de un punto, con audio, sin ninguna anotación. Única
%   preparación: marcar la pista a mano una vez por cámara.
% - Sale: por cada golpe, instante, jugador y tipo (derecha, revés, remate,
%   saque; agrupados por gesto: una volea de derecha es derecha).
% - El tipo no va solo: sin cuándo y quién no hay estadísticas por jugador.
%   Cadena; la pieza nueva es la última.
% - Listón: que funcione lo bastante bien para usarse de verdad, no una cifra
%   fijada de antemano.
% - Decorte et al., base del cuándo y el quién: no reinventar un método que
%   funcionaba; compararse solo para comprobar la reimplementación.
% - Visor: estadísticas de juego por jugador (golpes por tipo, total, mapa de
%   calor) y saltar al instante de cada golpe. Que sea útil.
% - Objetivos: detectar el instante; identificar al jugador; clasificar el tipo;
%   el visor.
% - No son objetivos: los datos etiquetados (medio: no existían; una línea en el
%   problema, el resto en Datos); la evaluación por torneos (metodología).
% - Fuera: marcador y resultado; bandeja, víbora, dejada; cámaras móviles o a ras
%   de pista; partidos completos (trabaja por puntos). Evaluado con
%   retransmisiones (únicas grabaciones anotadas, cumplen la condición de
%   cámara); otras alturas = trabajo futuro (PADELVIC, cuatro alturas).
% - Figura en el problema, de ejemplo y sin flechas: (a) fotograma crudo de un
%   punto; (b) lo que sale: línea de tiempo con todos los golpes y, sobre el
%   mismo fotograma, quién da el golpe de ese instante y de qué tipo. Punto de
%   VIGO, salida de los modelos entrenados sin ese torneo (chain_cv), golpe
%   correcto en instante, jugador y tipo.
% - Citas: Decorte et al.; CVSPORTS en el pie de la figura.

\section{Descripción del problema}
\label{sec:problema}

\section{Objetivo general}
\label{sec:objetivo-general}

\section{Objetivos específicos}
\label{sec:objetivos-especificos}

\section{Alcance}
\label{sec:alcance}
